\documentclass[10pt,a4paper]{amsart}
\usepackage[margin=19mm]{geometry}
\usepackage{amsmath,amssymb,mathtools}
\usepackage[scr=rsfs,cal=euler]{mathalfa}
\usepackage{xfrac}
\usepackage[unicode,hidelinks,bookmarks=false]{hyperref}
\newcommand{\ie}{i.e.\ }
\newcommand{\cf}{cf.\ }
\newcommand{\resp}{respectively\ }
\newcommand{\Subsection}{\subsection}
\providecommand{\nobreakdash}{\nobreak\hbox{-}}
\setlength{\emergencystretch}{2em}
\multlinegap0pt
\usepackage{amssymb}
\usepackage{mathtools}
\usepackage[shortlabels]{enumitem}
\usepackage{xcolor}
\usepackage{tikz}
\usepackage{cancel}
\newtheorem{proposition}{Proposition}
\newcommand{\Sone}{\mathbb{S}^1}
\newcommand{\Wone}{W_1}
\newcommand{\cA}{{\mathcal A}}
\newcommand{\cP}{{\mathcal P}}
\newcommand{\diam}{\operatorname{diam}}
\newcommand{\emp}{\mathsf e}
\newcommand{\wto}{\overset{*}{\rightharpoonup}}
\title{Low emergence in one-dimensional dynamics}
\author{Odylo Costa, Bruno Santiago}
\date{Source: arXiv:2608.14166v1 (CC BY 4.0)}
\begin{document}
\maketitle

\noindent\emph{Source and licence.} Low emergence in one-dimensional dynamics. Version arXiv:2608.14166v1 (CC BY 4.0), distributed by arXiv under the Creative Commons Attribution 4.0 International licence, http://creativecommons.org/licenses/by/4.0/, as recorded in the arXiv licence metadata for this exact version and shown on the abstract page https://arxiv.org/abs/2608.14166v1. Full licence notice: \texttt{licenses/arxiv-2608.14166-CC-BY-4.0.txt}.\par\smallskip

\noindent\textit{Editorial scope.} Counted unit: the article's proposition prop:iterate (source0.tex lines 1687--1692). The article's complete original proof of it is delivered with it (source0.tex lines 1694--1806, one \texttt{proof} environment of 113 lines). No mathematics is written, repaired or re-derived: the statement and the proof are the article's own lines, in the article's own notation, and no retained line is substituted or re-flowed.  No further source-local context is needed: every cross-reference of every retained line resolves inside the two delivered spans.  Nothing was omitted from any retained span, so the package carries no omission record.  No retained line contains a citation, so the wrapper carries no bibliography and no bibliography entry.  No retained line contains a figure environment, an image-inclusion command or drawing code, so no image is delivered and the package declares no asset dependency.  Editorial formatting only, in addition: an AMS \texttt{amsart} preamble that restates the article's own declarations and macros used by the retained lines, namely the environments \texttt{proposition} on separate counters, together with the standard packages those lines need.  The retained environments print in this document as Proposition 1 in order of appearance, whereas the article prints the same items with the numbering of its own full document.  The complete native source of the article as distributed in the arXiv e-print for this exact version is bundled unchanged as \texttt{sources/207641/source0.tex}.
\begin{proposition}\label{prop:iterate}
For every \(x\in\Sone\),
\[
\cA_f(x)=T_N(\cA_F(x)).
\]
\end{proposition}

\begin{proof}
For every \(q\ge1\),
\[
\emp_q^F(x)
=
\frac1q\sum_{k=0}^{q-1}\delta_{f^{kN}(x)}.
\]
Therefore,
\[
\begin{aligned}
T_N(\emp_q^F(x))
&=
\frac1N\sum_{j=0}^{N-1}
(f^j)_*
\left(
\frac1q\sum_{k=0}^{q-1}\delta_{f^{kN}(x)}
\right)
\\
&=
\frac1{qN}
\sum_{j=0}^{N-1}\sum_{k=0}^{q-1}
\delta_{f^{j+kN}(x)}
\\
&=
\frac1{qN}\sum_{m=0}^{qN-1}\delta_{f^m(x)}
=
\emp_{qN}^f(x).
\end{aligned}
\]
Thus,
\[
\emp_{qN}^f(x)=T_N(\emp_q^F(x)).
\]

Now write \(m=qN+r\), with \(0\le r<N\). Then
\[
\emp_{qN+r}^f(x)
=
\frac{qN}{qN+r}\,\emp_{qN}^f(x)
+
\frac1{qN+r}
\sum_{j=qN}^{qN+r-1}\delta_{f^j(x)},
\]
and hence
\[
\Wone(\emp_{qN+r}^f(x),\emp_{qN}^f(x))
\le
\frac{2r\,\diam(\Sone)}{qN+r}
\le
\frac{2N\,\diam(\Sone)}{qN+r}.
\]
In particular, this quantity tends to zero as \(q\to\infty\), uniformly
in \(x\) and \(0\le r<N\).

Let \(\eta\in\cA_F(x)\). There exists \(q_k\to\infty\) such that
\[
\emp_{q_k}^F(x)\wto\eta.
\]
Since \(T_N\) is continuous,
\[
\emp_{q_kN}^f(x)
=
T_N(\emp_{q_k}^F(x))
\wto T_N(\eta),
\]
and therefore
\[
T_N(\eta)\in\cA_f(x).
\]
Hence
\[
T_N(\cA_F(x))\subset\cA_f(x).
\]

Conversely, let \(\nu\in\cA_f(x)\). There exists \(m_k\to\infty\) such that
\[
\emp_{m_k}^f(x)\wto\nu.
\]
Write
\[
m_k=q_kN+r_k,
\qquad
0\le r_k<N.
\]
By the estimate above,
\[
\Wone(\emp_{m_k}^f(x),\emp_{q_kN}^f(x))\to0,
\]
and hence
\[
\emp_{q_kN}^f(x)\wto\nu.
\]
Since \(\cP(\Sone)\) is compact, after passing to a subsequence we may
assume that
\[
\emp_{q_k}^F(x)\wto\eta
\]
for some \(\eta\in\cA_F(x)\). Therefore,
\[
\nu
=
\lim_{k\to\infty}\emp_{q_kN}^f(x)
=
\lim_{k\to\infty}T_N(\emp_{q_k}^F(x))
=
T_N(\eta).
\]
Thus
\[
\cA_f(x)\subset T_N(\cA_F(x)),
\]
which proves the equality.
\end{proof}

\end{document}
